\phantomsection
\addcontentsline{toc}{chapter}{\bf ABSTRACT}

\begin{center}
    \large{\bf ABSTRACT}
\end{center}

Optimizing customer conversion paths is a central problem in computational marketing, where user interactions are modeled as a directed graph and the objective is to identify the highest-probability route from entry to conversion. As graph scale increases, exhaustive search becomes intractable. This work formulates the problem as a shortest-path query on a directed weighted graph whose edge weights are negative log-transition probabilities, reducing maximum-probability path search to standard shortest-path form.

A \textbf{Probability-Pruned Dijkstra} variant is proposed that introduces a threshold parameter $\tau$ to discard low-probability partial paths during traversal, reducing the effective edge set from $|E|$ to $|E'|\leq|E|$ with guaranteed convergence to the baseline as $\tau\to 0$. Five formal correctness theorems are established, including an all-or-nothing property: the pruned algorithm either returns the globally optimal path or reports no path at all.

Evaluation spans 3,000 configurations---2,160 synthetic (graphs up to 10,000 nodes) and 840 real-data rows from RetailRocket and RecSys~2015---under fixed and adaptive $\tau$ regimes. Wall-clock speedups range from $3.3\times$ to $738\times$ on synthetic data and reach $18{,}882\times$ on real-world graphs. An adaptive-$\tau$ regime raises path-found rates from near~0\% to over~80\%. Across all 706 path-found rows the optimality gap is exactly $0.0000\%$, confirming that threshold-based pruning provides practical acceleration without sacrificing solution quality.
